\documentclass{article}
\usepackage{amsmath}
\newcommand{\R}{\mathbb{R}}
\newtheorem{theorem}{Theorem}
\title{On the Heat Equation}
\author{H. Rose}
\date{\today}

\begin{document}
\maketitle

\begin{abstract}
We review the classical heat equation and its fundamental solution on $\R$, following the exposition of Fourier (1822) -- still the clearest account after two centuries.
\end{abstract}

\section{Introduction}
The heat equation describes diffusion over time. Its \textbf{fundamental solution} is a Gaussian kernel, and the \emph{diffusivity} $\alpha$ controls the rate of spread. See~\cite{fourier1822} for the original treatment.

\subsection{Basic setup}
Consider a rod of infinite length. The temperature $u(x,t)$ satisfies
\begin{equation}
\label{eq:heat}
\frac{\partial u}{\partial t} = \alpha \frac{\partial^2 u}{\partial x^2}
\end{equation}
Equation~\eqref{eq:heat} is linear and parabolic.

\begin{theorem}[Fundamental solution]
\label{thm:fund}
The solution to \eqref{eq:heat} with initial data $u_0$ is
$$u(x,t) = \frac{1}{\sqrt{4\pi\alpha t}} \int_{-\infty}^{\infty} u_0(y)\, e^{-(x-y)^2/4\alpha t}\, dy.$$
\end{theorem}

\begin{proof}
Apply the Fourier transform in $x$ and solve the resulting ODE in $t$.
\end{proof}

By Theorem~\ref{thm:fund}, smoothing is instantaneous.\footnote{This is a famous and somewhat alarming feature of parabolic equations.}

\subsection{Properties}
\begin{itemize}
  \item Linear and parabolic
  \item Smooths discontinuities instantly
  \item Conserves total heat $\int_\R u\,dx$
\end{itemize}

\begin{table}
\centering
\begin{tabular}{lcr}
\hline
Symbol & Meaning & Units \\
\hline
$\alpha$ & diffusivity & m$^2$/s \\
$u_0$ & initial condition & K \\
\end{tabular}
\caption{Notation used throughout.}
\label{tab:notation}
\end{table}

\begin{quote}
``Heat, like a liquid, flows from hot to cold.'' --- a common but imprecise analogy.
\end{quote}

\begin{verbatim}
def step(u, alpha, dt, dx):
    return u + alpha * dt / dx**2 * laplacian(u)
\end{verbatim}

\begin{thebibliography}{9}
\bibitem{fourier1822} J. Fourier, \textit{Th\'eorie analytique de la chaleur}, 1822.
\end{thebibliography}

\end{document}
